% !TEX program = xelatex
\documentclass[final]{beamer}

\usepackage[size=a0,orientation=landscape,scale=1.00]{beamerposter}
\usepackage{fontspec}
\usepackage{amsmath,amssymb,mathtools}
\usepackage{booktabs,array}
\usepackage{ragged2e}
\usepackage{tikz}
\usetikzlibrary{arrows.meta,calc,positioning,shapes.geometric}
\usepackage[most]{tcolorbox}

\setsansfont{TeX Gyre Heros}
\setmonofont{TeX Gyre Cursor}

\definecolor{Navy}{HTML}{0B1F33}
\definecolor{NavyTwo}{HTML}{163A53}
\definecolor{Teal}{HTML}{0E8F8B}
\definecolor{TealLight}{HTML}{DDF4F1}
\definecolor{Amber}{HTML}{E8A72E}
\definecolor{AmberLight}{HTML}{FFF2D5}
\definecolor{Ink}{HTML}{16212B}
\definecolor{Muted}{HTML}{536573}
\definecolor{Paper}{HTML}{F6F8F9}
\definecolor{White}{HTML}{FFFFFF}
\definecolor{Rule}{HTML}{CBD5DB}
\definecolor{Red}{HTML}{B53A4C}

\setbeamercolor{background canvas}{bg=Paper}
\setbeamercolor{normal text}{fg=Ink,bg=Paper}
\setbeamertemplate{navigation symbols}{}
\setbeamertemplate{headline}{}
\setbeamertemplate{footline}{}
\setbeamersize{text margin left=2.5cm,text margin right=2.5cm}
\hypersetup{
  pdftitle={Orden causal y cuentas uniformes no fijan la distancia por cadenas},
  pdfauthor={Ignacio Martín Gandul},
  pdfsubject={WP7: contraejemplo cuantitativo a F1+F2 implica F3 en dimensión 1+1},
  pdfkeywords={causal sets, Madsen, Muller, de Sitter, discrepancy, longest chain, F1, F2, F3}
}

\newtcolorbox{posterbox}[2][]{
  enhanced,
  colback=White,
  colframe=Rule,
  boxrule=0.9pt,
  arc=3mm,
  left=8mm,right=8mm,top=6mm,bottom=7mm,
  before skip=0mm,after skip=9mm,
  title={#2},
  colbacktitle=Navy,
  coltitle=White,
  fonttitle=\bfseries\fontsize{28}{32}\selectfont,
  attach boxed title to top left={xshift=0mm,yshift=-1mm},
  boxed title style={arc=2.5mm,boxrule=0pt,left=7mm,right=7mm,top=3mm,bottom=3mm},
  #1
}

\newtcolorbox{resultbox}[1][]{
  enhanced,
  colback=TealLight,
  colframe=Teal,
  boxrule=2.2pt,
  arc=3mm,
  left=10mm,right=10mm,top=8mm,bottom=8mm,
  before skip=0mm,after skip=9mm,
  #1
}

\newtcolorbox{scopebox}[1][]{
  enhanced,
  colback=AmberLight,
  colframe=Amber,
  boxrule=1.4pt,
  arc=3mm,
  left=8mm,right=8mm,top=6mm,bottom=6mm,
  before skip=0mm,after skip=9mm,
  #1
}

\newcommand{\badge}[2]{%
  \tikz[baseline=(X.base)]\node[
    rounded corners=2.4mm,
    inner xsep=5.5mm,
    inner ysep=2.6mm,
    fill=#1,
    text=White,
    font=\bfseries\fontsize{15}{18}\selectfont
  ] (X) {#2};%
}

\newcommand{\key}[1]{\textcolor{Teal}{\bfseries #1}}
\newcommand{\warn}[1]{\textcolor{Red}{\bfseries #1}}
\newcommand{\tightdisplay}[1]{\vspace{-2mm}\[#1\]\vspace{-3mm}}

\begin{document}
\begin{frame}[t]

% Header
\begin{tcolorbox}[
  enhanced,
  colback=Navy,
  colframe=Navy,
  boxrule=0pt,
  arc=0mm,
  left=16mm,right=16mm,top=10mm,bottom=10mm,
  width=\textwidth
]
\begin{columns}[T,totalwidth=\linewidth]
  \begin{column}{0.82\linewidth}
    {\fontsize{60}{64}\selectfont\bfseries\color{White}
      Orden causal y cuentas uniformes no fijan la distancia por cadenas\par}
    \vspace{4mm}
    {\fontsize{31}{36}\selectfont\color{TealLight}
      Un contraejemplo cuantitativo a \(\mathrm{F1}+\mathrm{F2}\Rightarrow\mathrm{F3}\)
      en dimensión \(1+1\)\par}
    \vspace{5mm}
    {\fontsize{19}{23}\selectfont\color{White}
      Ignacio Martín Gandul \quad\textcolor{Teal}{\textbullet}\quad
      Programa \texttt{nachocausal} \quad\textcolor{Teal}{\textbullet}\quad
      WP7 · 11 agosto 2026}
  \end{column}
  \begin{column}{0.18\linewidth}
    \raggedleft
    \begin{tikzpicture}
      \fill[Teal] (0,0) circle (2.15cm);
      \draw[White,line width=2.5pt] (0,0) circle (2.15cm);
      \node[White,font=\bfseries\fontsize{38}{40}\selectfont] at (0,0.15) {WP7};
      \node[White,font=\fontsize{14}{16}\selectfont] at (0,-0.75) {\(1+1\)D};
    \end{tikzpicture}
  \end{column}
\end{columns}
\end{tcolorbox}

\vspace{8mm}
\fontsize{22.5}{27.5}\selectfont

\begin{columns}[T,totalwidth=\textwidth]

% ------------------------------------------------------------------
% COLUMN 1
% ------------------------------------------------------------------
\begin{column}{0.318\textwidth}

\begin{resultbox}
  {\fontsize{30}{34}\selectfont\bfseries\color{Navy}Resultado central}\par
  \vspace{4mm}
  Bajo la lectura regional explícita de la Remark 5.4 de Madsen y cualquiera de las dos
  convenciones estándar auditadas para la norma de curvatura, existen configuraciones finitas
  de alta densidad en de Sitter planar \(1+1\) que:
  \vspace{3mm}
  \begin{itemize}
    \setlength\itemsep{2mm}
    \item preservan exactamente el orden causal (F1);
    \item satisfacen F2 con una constante fija, hasta la escala mínima exacta;
    \item violan la desigualdad completa F3 por un margen constante.
  \end{itemize}
  \vspace{2mm}
  \centering
  \badge{Teal}{F1 + F2 \(\not\Rightarrow\) F3}
\end{resultbox}

\begin{posterbox}{La pregunta de Madsen}
  Para un embedding de un causal set en un espaciotiempo, Madsen separa tres exigencias:
  \vspace{3mm}

  \begin{tabular}{@{}>{\bfseries\color{Navy}}p{0.10\linewidth}p{0.84\linewidth}@{}}
    F1 & El orden discreto coincide exactamente con el orden causal inducido.\\[2mm]
    F2 & Las cuentas en todos los diamantes admisibles aproximan el volumen con tolerancia
         cuantitativa uniforme.\\[2mm]
    F3 & La longitud de cadena aproxima el tiempo propio con la tasa prescrita.\\
  \end{tabular}

  \vspace{5mm}
  \begin{center}
  \begin{tikzpicture}[node distance=7mm]
    \node[draw=Teal,fill=TealLight,rounded corners=2mm,minimum width=3.3cm,
          minimum height=1.4cm,font=\bfseries] (f1) {F1};
    \node[draw=Teal,fill=TealLight,rounded corners=2mm,minimum width=3.3cm,
          minimum height=1.4cm,right=of f1,font=\bfseries] (f2) {F2};
    \node[draw=Red,fill=White,rounded corners=2mm,minimum width=3.3cm,
          minimum height=1.4cm,right=2.2cm of f2,font=\bfseries] (f3) {F3};
    \node[font=\bfseries\Large] at ($(f1)!0.5!(f2)$) {+};
    \draw[-{Latex[length=6mm]},line width=2pt,Navy] (f2) -- (f3);
    \draw[Red,line width=3pt] ($(f2.east)+(0.72,0.58)$) -- ($(f2.east)+(1.48,-0.58)$);
    \draw[Red,line width=3pt] ($(f2.east)+(0.72,-0.58)$) -- ($(f2.east)+(1.48,0.58)$);
    \node[below=2mm of f3,font=\fontsize{15}{18}\selectfont,text=Red] {no es redundante};
  \end{tikzpicture}
  \end{center}
\end{posterbox}

\begin{posterbox}{La tolerancia que sí se conserva}
  Para \(\mathcal N_\rho=\rho V_K\), la configuración final \(P_\rho\) satisface,
  uniformemente en todo diamante F2-admisible \(D\),
  \tightdisplay{
    \bigl|\#(P_\rho\cap D)-\rho\operatorname{Vol}_g(D)\bigr|
    \le K_2\sqrt{\rho\operatorname{Vol}_g(D)\log\mathcal N_\rho},
  }
  donde \(K_2\) es fija e independiente de \(\rho\), de \(D\) y de su posición.

  \vspace{4mm}
  La escala mínima mantiene
  \tightdisplay{
    \rho\operatorname{Vol}_g(D)
    \ge \frac{\kappa_\beta}{c_*^2}\log^2\mathcal N_\rho.
  }
  Ésta es la tolerancia cuantitativa exacta relevante: no una noción informal de
  “uniformidad local”.
\end{posterbox}

\begin{posterbox}{Mapa auditable de la demostración}
  La prueba está separada en obligaciones que pueden verificarse independientemente:
  \vspace{3mm}

  \renewcommand{\arraystretch}{1.28}
  \begin{tabular}{@{}>{\bfseries\color{Navy}}p{0.105\linewidth}p{0.65\linewidth}>{\raggedleft\arraybackslash}p{0.19\linewidth}@{}}
    P1 & Fondo de Hammersley con discrepancia \(O(\!\log N)\) & \textcolor{Teal}{\bfseries PROVED}\\
    P2 & Control de todos los rectángulos, incluidos los alargados & \textcolor{Teal}{\bfseries PROVED}\\
    P3 & Cadena de \(\Theta(\sqrt N)\) compatible con F2 & \textcolor{Teal}{\bfseries PROVED}\\
    P4 & Violación de F3 por factor constante & \textcolor{Teal}{\bfseries PROVED}\\
    P5 & Transferencia al parche lorentziano \(1+1\) & \textcolor{Amber}{\bfseries SCOPED}\\
  \end{tabular}

  \vspace{4mm}
  Dentro de P5.2: O1 y O4 están probadas; O2 pasa bajo ambas normas estándar; O3 conserva
  alcance regional; O5 queda \texttt{PROVED\_MODULO\_O3\_SCOPE}.
\end{posterbox}

\end{column}

\hfill

% ------------------------------------------------------------------
% COLUMN 2
% ------------------------------------------------------------------
\begin{column}{0.345\textwidth}

\begin{posterbox}{Mecanismo: fondo cuasiuniforme + cadena plantada}
  \begin{columns}[T,totalwidth=\linewidth]
    \begin{column}{0.48\linewidth}
      \begin{center}
      \begin{tikzpicture}[x=0.52cm,y=0.52cm]
        \fill[Paper] (0,0) rectangle (14,14);
        \draw[Navy,line width=1.8pt] (0,0) rectangle (14,14);
        \draw[Rule,line width=0.8pt,step=2] (0,0) grid (14,14);
        \fill[TealLight,opacity=.75] (2.2,1.6) rectangle (11.8,10.6);
        \draw[Teal,dashed,line width=1.8pt] (2.2,1.6) rectangle (11.8,10.6);
        % Schematic quasiuniform background
        \foreach \x/\y in {
          .7/8.3,1.4/3.7,2.1/12.5,2.8/6.1,3.5/1.0,4.2/9.8,4.9/4.4,
          5.6/13.1,6.3/7.2,7.0/2.4,7.7/11.4,8.4/5.2,9.1/13.7,9.8/8.8,
          10.5/3.1,11.2/10.2,11.9/6.5,12.6/1.8,13.3/12.0,1.0/10.9,
          3.1/8.0,5.0/1.9,6.8/12.6,8.8/9.9,10.8/5.8,12.9/7.7}
          \fill[NavyTwo] (\x,\y) circle (0.13);
        % Planted chain
        \draw[Amber,line width=3.0pt] (3.1,3.0) -- (10.9,10.8);
        \foreach \t in {0,...,12}{
          \pgfmathsetmacro{\u}{3.1+0.65*\t}
          \pgfmathsetmacro{\v}{3.0+0.65*\t}
          \fill[Amber] (\u,\v) circle (0.19);
          \draw[White,line width=.6pt] (\u,\v) circle (0.19);
        }
        \node[anchor=west,font=\fontsize{13}{15}\selectfont,text=Navy] at (0,14.6)
          {Esquema en coordenadas nulas};
        \node[anchor=west,font=\fontsize{12}{14}\selectfont,text=Teal] at (2.3,1.0)
          {diamante \(\leftrightarrow\) rectángulo};
        \node[rotate=45,font=\bfseries\fontsize{12}{14}\selectfont,text=Amber] at (9.8,11.7)
          {cadena \(\Gamma_\rho\)};
      \end{tikzpicture}
      \end{center}
    \end{column}
    \begin{column}{0.49\linewidth}
      En coordenadas nulas de \(1+1\):
      \vspace{2mm}
      \begin{itemize}
        \setlength\itemsep{2.1mm}
        \item causalidad \(\leftrightarrow\) orden producto;
        \item diamantes \(\leftrightarrow\) rectángulos;
        \item cadenas \(\leftrightarrow\) subsecuencias crecientes.
      \end{itemize}
      \vspace{3mm}
      Se parte de un fondo \(\Pi_\rho\) que ya satisface F1--F2 y se añade, sobre una
      geodésica temporal fija de longitud \(\tau_0\),
      \tightdisplay{
        \Gamma_\rho,\qquad
        k_\rho=\Bigl\lceil A\sqrt{m_2\rho}\,\tau_0\Bigr\rceil.
      }
      La masa añadida es sólo \(O(\!\rho^{-1/2})\), pero está perfectamente alineada en el orden.
    \end{column}
  \end{columns}
\end{posterbox}

\begin{posterbox}{Por qué F2 sobrevive}
  La intersección de la geodésica plantada con cualquier diamante causalmente convexo es un
  intervalo. En de Sitter planar, una comparación uniforme volumen--traza da
  \tightdisplay{
    \#(\Gamma_\rho\cap D)
    \le A\sqrt{\frac{m_2}{\kappa_\beta}}
       \sqrt{\rho\operatorname{Vol}_g(D)}+2.
  }
  Al normalizar por el presupuesto de F2, el coste adicional es
  \tightdisplay{
    \Xi_A^{\mathrm{dS}}(\rho)=
    \frac{A\sqrt{m_2/\kappa_\beta}}{\sqrt{\log\mathcal N_\rho}}
    +\frac{2c_*}{\sqrt{\kappa_\beta}\,\log^{3/2}\mathcal N_\rho}
    \longrightarrow 0.
  }
  Por tanto, para todo \(K_2>K_{\rm bg}\), F2 vuelve a cumplirse con \(K_2\) fija a partir de
  densidad suficientemente grande.
\end{posterbox}

\begin{posterbox}{Por qué F3 falla por margen constante}
  Para los extremos plantados \(x_\rho\prec y_\rho\), la cadena impone
  \tightdisplay{
    \frac{\ell_{P_\rho}(x_\rho,y_\rho)}{\sqrt{m_2\rho}}-\tau_0
    \ge (A-1)\tau_0-o(1).
  }
  F3 sólo permitiría
  \tightdisplay{
    C_2\!\left(\varepsilon^2+
    \frac{\log^{3/2}\mathcal N_\rho}{\mathcal N_\rho^{1/4}}\right)\tau_0.
  }
  Elegir \(A>1+C_2\varepsilon^2\) deja un margen fijo
  \(\mu=A-1-C_2\varepsilon^2>0\). Para \(\rho\) grande, la violación es al menos
  \(\mu\tau_0/3\): falla la \key{desigualdad completa}, no sólo una tasa secundaria.
\end{posterbox}

\begin{scopebox}
  \centering
  \badge{NavyTwo}{O5: PROVED\_MODULO\_O3\_SCOPE}
  \hspace{6mm}
  \badge{Teal}{P5.2: PASS\_WITH\_SCOPE}
\end{scopebox}

\begin{posterbox}{La ventana de escalas que hace posible el contraejemplo}
  \begin{center}
  \begin{tikzpicture}[x=1cm,y=1cm]
    \draw[-{Latex[length=5mm]},line width=2pt,Navy] (0,0) -- (15.8,0);
    \node[below,font=\fontsize{13}{15}\selectfont,text=Muted] at (15.1,-0.15) {densidad \(\rho\)};

    \fill[NavyTwo] (1.0,0) circle (.17);
    \node[above=3mm,align=center,font=\fontsize{14}{17}\selectfont,text=Navy] at (3.0,0)
      {fondo\\\(\mathcal N_\rho=\Theta(\rho)\)};

    \fill[Amber] (7.8,0) circle (.20);
    \node[above=3mm,align=center,font=\bfseries\fontsize{14}{17}\selectfont,text=Amber] at (7.8,0)
      {cadena\\\(k_\rho=\Theta(\sqrt\rho)\)};

    \fill[Teal] (14.0,0) circle (.20);
    \node[above=3mm,align=center,font=\fontsize{14}{17}\selectfont,text=Teal] at (13.1,0)
      {fracción añadida\\\(\Theta(\rho^{-1/2})\to0\)};

    \draw[Teal,line width=1.6pt,{Latex[length=4mm]}-] (4.1,-1.25) -- (1.0,-1.25);
    \draw[Teal,line width=1.6pt,-{Latex[length=4mm]}] (4.1,-1.25) -- (7.2,-1.25);
    \node[below=1mm,align=center,font=\fontsize{13}{16}\selectfont,text=Teal] at (4.1,-1.25)
      {coste F2 normalizado\\\(\Theta((\log\mathcal N_\rho)^{-1/2})\to0\)};

    \draw[Red,line width=1.6pt,{Latex[length=4mm]}-] (11.1,-1.25) -- (8.1,-1.25);
    \draw[Red,line width=1.6pt,-{Latex[length=4mm]}] (11.1,-1.25) -- (14.1,-1.25);
    \node[below=1mm,align=center,font=\bfseries\fontsize{13}{16}\selectfont,text=Red] at (11.1,-1.25)
      {sesgo F3 relativo\\\(\Theta(1)\)};
  \end{tikzpicture}
  \end{center}

  \vspace{2mm}
  La misma perturbación es asintóticamente invisible para el presupuesto de cuentas y
  macroscópicamente visible para la altura de la cadena. Éste es el núcleo cuantitativo de WP7.
\end{posterbox}

\end{column}

\hfill

% ------------------------------------------------------------------
% COLUMN 3
% ------------------------------------------------------------------
\begin{column}{0.318\textwidth}

\begin{posterbox}{Müller: el precursor causal negativo más cercano}
  Müller demuestra que, para \(K\) fijo, existen dos slabs de Cauchy de volumen uno, con las
  mismas fronteras, distancia \(d^-\) arbitrariamente grande y leyes de órdenes de \(K\) muestras
  arbitrariamente próximas. El resultado persiste bajo su “Planck-scale uniformness” por sí sola.

  \vspace{5mm}
  \renewcommand{\arraystretch}{1.23}
  \begin{tabular}{@{}>{\color{Navy}\bfseries}p{0.43\linewidth}p{0.51\linewidth}@{}}
    Müller & WP7\\
    \midrule
    \(K\) fijo & \(\rho\to\infty\)\\
    Dos geometrías y leyes próximas & Una configuración y un embedding\\
    Distancia global \(d^-\) & Cadena que viola F3\\
    Error absoluto \(Ks\) & Tolerancia exacta de Madsen\\
    Uniformidad sin tasa comparable & Escala mínima y cuantificadores explícitos\\
  \end{tabular}

  \vspace{5mm}
  \begin{center}
    \badge{Amber}{CLOSEST CAUSAL NEGATIVE PRECURSOR}
    \hspace{3mm}
    \badge{NavyTwo}{DOES NOT SUBSUME}
  \end{center}
\end{posterbox}

\begin{posterbox}{Realización geométrica: de Sitter planar}
  El candidato único auditado es
  \tightdisplay{
    M_\ell=(-\infty,0)_\eta\times\mathbb R_x,
    \qquad
    g=\frac{\ell^2}{\eta^2}\bigl(-d\eta^2+dx^2\bigr).
  }
  Proporciona hiperbolicidad global, una escala de curvatura finita bajo ambas normas estándar,
  una región precompacta con diamante testigo profundo y una cota volumen--traza uniforme.

  \vspace{4mm}
  \warn{Dos calificadores deben acompañar siempre a “de Sitter planar”:}
  \vspace{3mm}
  \begin{center}
  \begin{tcolorbox}[colback=Paper,colframe=Amber,boxrule=1pt,arc=2mm,
                    left=4mm,right=4mm,top=3mm,bottom=3mm,width=.94\linewidth]
    \centering\bfseries\ttfamily\fontsize{15}{19}\selectfont
    CURVATURE\_NORM\_UNDEFINED\\
    BOTH\_STANDARD\_READINGS\_PASS\\[2mm]
    REGIONAL\_SCOPE\_NOT\_LITERAL\\
    EXPLICIT\_REGIONAL\_READING\_ADOPTED
  \end{tcolorbox}
  \end{center}
\end{posterbox}

\begin{posterbox}{Qué es nuevo y qué no se afirma}
  \key{Separación específica.} WP7 combina en una misma sucesión F1 exacta, F2 con la tolerancia
  cuantitativa de Madsen y una cadena concreta que viola F3.

  \vspace{4mm}
  \key{Precedentes obligatorios.}
  \begin{itemize}
    \setlength\itemsep{1.8mm}
    \item \textbf{Henson}: cambiar pocas relaciones causales puede destruir embebibilidad;
          no habla de añadir elementos.
    \item \textbf{Braun}: “order + number = geometry” para leyes de todo tamaño y \(d\ge3\);
          no cubre \(1+1\) ni F2\(\Rightarrow\)F3.
    \item \textbf{Dick--Pillichshammer}: el Teorema 3.46, p. impresa 100 (PDF 99), da exactamente
          \(ND_N^*\le\log_2N+4\) al especializar a \(s=2,b_1=2\).
  \end{itemize}

  \vspace{3mm}
  \warn{No se afirma} prioridad absoluta, “independencia lógica” sin calificativos, falsedad del
  teorema de Madsen, extensión a \(d\ge3\) ni reconstrucción métrica general.
\end{posterbox}

\begin{posterbox}[colback=Navy,colframe=Navy,coltitle=White]{Formulación pública segura}
  \color{White}
  \textbf{We answer negatively the question left open in Madsen's Remark/footnote: under the
  stated regional interpretation in \(1+1\) dimensions, F1 and F2 do not imply F3.}

  \vspace{4mm}
  \color{TealLight}
  Müller is the closest causal negative precursor; WP7 isolates a distinct, more specific
  separation using F1--F3 and Madsen's exact quantitative tolerance.
\end{posterbox}

\end{column}
\end{columns}

\vfill

% Footer
\begin{tcolorbox}[
  enhanced,colback=NavyTwo,colframe=NavyTwo,boxrule=0pt,arc=0mm,
  left=10mm,right=10mm,top=5mm,bottom=5mm,width=\textwidth
]
{\fontsize{14.5}{18}\selectfont\color{White}
\textbf{Fuentes clave.}
N. Madsen, arXiv:2607.05840v1 (2026), Defs. 2.2, 2.6 y Remarks 4.7, 5.4;
O. Müller, arXiv:2503.01719v2 (2025), Thm. 2 y pp. 3--4;
J. Dick y F. Pillichshammer, \emph{Digital Nets and Sequences} (2010), Thm. 3.46, p. 100;
J. Henson, gr-qc/0601121, §1.3;
M. Braun, \emph{Class. Quantum Grav.} 43 (2026), 045015.
\hfill
\textbf{Estado:} verificación literal cerrada · sin claim de prioridad absoluta.}
\end{tcolorbox}

\end{frame}
\end{document}
